\chapter{Geometria Computacional 2D}

%==============================================================================
\section{Template 2D}
%==============================================================================

\subsection{Padrão do Problema}
Problemas com pontos, vetores, segmentos, orientação, distância e interseção no plano.

\subsection{Quando usar}
\begin{itemize}[nosep]
  \item O enunciado trabalha com coordenadas no plano.
  \item É preciso decidir se três pontos fazem curva à esquerda, à direita ou são colineares.
  \item Há consultas sobre segmentos, distâncias ou interseções.
  \item Comparações com \texttt{double} precisam de tolerância.
\end{itemize}

\subsection{Explicação do Algoritmo}
O template centraliza operações geométricas básicas: produto escalar, produto vetorial, orientação \texttt{ccw}, teste de ponto em segmento, interseção de segmentos e distância de ponto a segmento.

\subsection{Complexidade Temporal}
$O(1)$ por operação geométrica básica.

\subsection{Exemplo de Enunciado}
\begin{exemplodiabox}
\textbf{Enunciado:} Dados dois segmentos $AB$ e $CD$, determine se eles se intersectam.

\textbf{Exemplo:} Se os segmentos cruzam em algum ponto ou compartilham extremidade, a resposta é sim.
\end{exemplodiabox}

\subsection{Código}
\incluircppcaption{algoritmos_avulsos/geometria/template_2d.cpp}{Template básico de geometria 2D}

%==============================================================================
\section{Ângulo entre vetores}
%==============================================================================

\subsection{Padrão do Problema}
Calcular o ângulo formado por três pontos ou comparar direção entre dois vetores.

\subsection{Quando usar}
\begin{itemize}[nosep]
  \item O problema pede ângulo entre segmentos.
  \item É necessário saber a rotação de um vetor para outro.
  \item Produto escalar e produto vetorial descrevem a relação entre direções.
  \item A resposta pode ser dada em radianos ou convertida para graus.
\end{itemize}

\subsection{Explicação do Algoritmo}
Para o ângulo $PQR$, forme os vetores $\overrightarrow{QP}$ e $\overrightarrow{QR}$. Use \texttt{atan2(cross, dot)} para obter o ângulo orientado com estabilidade numérica.

\subsection{Complexidade Temporal}
$O(1)$.

\subsection{Exemplo de Enunciado}
\begin{exemplodiabox}
\textbf{Enunciado:} Dados três pontos $P$, $Q$ e $R$, calcule o ângulo formado no vértice $Q$.
\end{exemplodiabox}

\subsection{Código}
\incluircppcaption{algoritmos_avulsos/geometria/angulo_vetores.cpp}{Ângulo entre vetores}

%==============================================================================
\section{Área de polígono}
%==============================================================================

\subsection{Padrão do Problema}
Calcular a área de um polígono simples a partir dos vértices em ordem.

\subsection{Quando usar}
\begin{itemize}[nosep]
  \item O enunciado dá os vértices de um polígono no plano.
  \item É preciso calcular área sem triangular manualmente.
  \item Os vértices estão em ordem horária ou anti-horária.
  \item O polígono pode ser convexo ou côncavo, desde que simples.
\end{itemize}

\subsection{Explicação do Algoritmo}
A fórmula de Shoelace soma os produtos vetoriais de vértices consecutivos. O valor absoluto da soma dividido por 2 é a área do polígono.

\subsection{Complexidade Temporal}
$O(n)$.

\subsection{Exemplo de Enunciado}
\begin{exemplodiabox}
\textbf{Enunciado:} Dado um terreno representado por $n$ vértices em ordem, calcule sua área.

\textbf{Exemplo:} Um retângulo com lados 3 e 4 tem área 12.
\end{exemplodiabox}

\subsection{Código}
\incluircppcaption{algoritmos_avulsos/geometria/area_poligono.cpp}{Área de polígono}

%==============================================================================
\section{Fecho convexo}
%==============================================================================

\subsection{Padrão do Problema}
Encontrar o menor polígono convexo que contém todos os pontos dados.

\subsection{Quando usar}
\begin{itemize}[nosep]
  \item O problema pede a borda externa de um conjunto de pontos.
  \item É preciso descartar pontos internos.
  \item A solução depende de ordenar pontos e manter apenas curvas válidas.
  \item Aplicações envolvem casco, cerca, contorno ou envoltória convexa.
\end{itemize}

\subsection{Explicação do Algoritmo}
O algoritmo de monotonic chain ordena os pontos e constrói duas partes do fecho: inferior e superior. Enquanto o último ponto cria uma curva inválida, ele é removido da pilha.

\subsection{Complexidade Temporal}
$O(n \log n)$ por causa da ordenação.

\subsection{Exemplo de Enunciado}
\begin{exemplodiabox}
\textbf{Enunciado:} Dados pontos de postes em um campo, encontre o contorno convexo mínimo que cerca todos eles.
\end{exemplodiabox}

\subsection{Código}
\incluircppcaption{algoritmos_avulsos/geometria/convex_hull.cpp}{Fecho convexo}

%==============================================================================
\section{Par de pontos mais próximo}
%==============================================================================

\subsection{Padrão do Problema}
Encontrar a menor distância entre dois pontos em um conjunto grande.

\subsection{Quando usar}
\begin{itemize}[nosep]
  \item Há muitos pontos e testar todos os pares seria lento.
  \item O problema pede menor distância euclidiana.
  \item A solução $O(n^2)$ não cabe.
  \item Os pontos podem ser ordenados por coordenada.
\end{itemize}

\subsection{Explicação do Algoritmo}
Divida os pontos pela coordenada $x$, resolva recursivamente cada metade e depois teste apenas pontos próximos da linha central. Ordenar por $y$ na combinação limita a quantidade de comparações locais.

\subsection{Complexidade Temporal}
$O(n \log n)$.

\subsection{Exemplo de Enunciado}
\begin{exemplodiabox}
\textbf{Enunciado:} Dadas as coordenadas de $n$ antenas, encontre a menor distância entre duas antenas distintas.
\end{exemplodiabox}

\subsection{Código}
\incluircppcaption{algoritmos_avulsos/geometria/closest_pair.cpp}{Par de pontos mais próximo}
